\section{Analyse des résultats sur AMD}

\begin{frame}{Identification des sections critiques (Option -R1)}
    \begin{columns}[T]
        \begin{column}{0.45\textwidth}
            \begin{itemize}
                \fontsize{8}{9}\selectfont
                \item Localisation des fonctions consommant le plus de temps CPU.
                \item \textbf{resid} est la routine dominante, représentant près de la moitié de la charge globale.
            \end{itemize}
            \vspace{0.1cm}
            \begin{table}
                \fontsize{6}{7}\selectfont
                \centering
                \begin{tabular}{|l|c|}
                    \hline
                    \textbf{Fonction} & \textbf{\% Temps CPU} \\ \hline
                    \texttt{resid}    & 47,85 \% \\ \hline
                    \texttt{psinv}    & 19,36 \% \\ \hline
                    \texttt{interp}   & 9,50 \%  \\ \hline
                    \texttt{rprj3}    & 6,69 \%  \\ \hline
                \end{tabular}
            \end{table}
        \end{column}
        \begin{column}{0.55\textwidth}
            \vspace{-0.4cm}
            \inputImage{figures/top_functions_amd_r1.png}{Analyse -R1 : Routines dominantes (AMD)}{0.95\textwidth}{!}{fig:r1_amd}
        \end{column}
    \end{columns}
\end{frame}

\begin{frame}{Fiabilité et Stabilité des mesures (Option -S1)}
    \begin{itemize}
        \fontsize{8}{9}\selectfont
        \item Validation de la robustesse des données via \textbf{10 répétitions}.
        \item Élimination des bruits de mesure liés au système.
    \end{itemize}
    
    \inputImage{figures/stability_analysis_amd_s1.png}{Analyse -S1 : Stabilité des mesures (AMD)}{0.85\textwidth}{0.45\textheight}{fig:s1_amd}
    
    \begin{block}{\footnotesize Constat de stabilité}
        \fontsize{7}{8}\selectfont
        Temps moyen stabilisé à \textbf{13,11 s}. La faible variance confirme la fiabilité de l'environnement de test pour les comparaisons de performance.
    \end{block}
\end{frame}

\begin{frame}{Étude de la montée en charge (Option -WS)}
    \begin{itemize}
        \fontsize{8}{9}\selectfont
        \item Évaluation de l'efficacité de l'implémentation parallèle de 1 à 8 cœurs.
    \end{itemize}
    
    \inputImage{figures/scalability_speedup_amd_ws.png}{Analyse -WS : Scalabilité (AMD)}{0.85\textwidth}{0.4\textheight}{fig:ws_amd}
    
    \begin{columns}[T]
        \begin{column}{0.48\textwidth}
            \begin{block}{\footnotesize Performance Parallèle}
                \fontsize{7}{8}\selectfont
                Accélération notable : passage de \textbf{30,27 s} à \textbf{10,73 s}.
            \end{block}
        \end{column}
        \begin{column}{0.48\textwidth}
            \begin{alertblock}{\footnotesize Limites de scalabilité}
                \fontsize{7}{8}\selectfont
                Apparition d'un \textit{Scalability Gap} de \textbf{2,84}. Le profil est limité par la bande passante mémoire (\textbf{Memory Bound}).
            \end{alertblock}
        \end{column}
    \end{columns}
\end{frame}

\begin{frame}{Duel de Paradigmes : Fortran 90 vs C++}
    \vspace{-0.2cm}
    \begin{itemize}
        \fontsize{8}{9}\selectfont
        \item Comparaison des deux implémentations avec le compilateur AOCC.
    \end{itemize}
    
    \inputImage{figures/comparison_duel_report_amd.png}{Comparaison de performance : Fortran vs C++}{0.85\textwidth}{0.4\textheight}{fig:duel_amd}
    
    \begin{columns}[T]
        \begin{column}{0.48\textwidth}
            \begin{exampleblock}{\footnotesize Fortran (Référence)}
                \fontsize{7}{8}\selectfont
                Score de compilation : \textbf{100 \%}. Exploitation maximale des unités de calcul vectorielles.
            \end{exampleblock}
        \end{column}
        \begin{column}{0.48\textwidth}
            \begin{alertblock}{\footnotesize C++ (NPB-CPP)}
                \fontsize{7}{8}\selectfont
                Score de compilation : \textbf{16,7 \%}. L'aliasing des pointeurs empêche la vectorisation automatique.
            \end{alertblock}
        \end{column}
    \end{columns}
\end{frame}